\section{Marco de options e implementación de SMDP}

\subsection{Formulación Semi-MDP y criterios de modelado}

El marco de options extiende el MDP a un Semi-MDP, lo que permite que la policy elija macro-actions que duran varios pasos. Denotamos con $\tau_\omega$ la longitud efectiva de ejecución de una option; el reward esperado se modela como $\mathbb{E}\left[\sum_{t=0}^{\tau_\omega} \gamma^t r_t\right]$, y el discounted factor conserva la convergencia.

\begin{importantbox}{La terna del diseño de una option}
\begin{itemize}
    \item $\mathcal{I}_\omega$: en qué estados puede activarse la option;
    \item $\pi_\omega$: la política interna se encarga de las acciones de varios pasos;
    \item $\beta_\omega$: la función de terminación controla el horizon de la option.
\end{itemize}
\end{importantbox}

\subsection{Estimación de valor y ecuación de Bellman}

La función value de una option se expresa como \(\mathrm{Q}(s, \omega) = \mathbb{E}\left[\sum_{t=0}^{\tau_\omega} \gamma^t r_t + \gamma^{\tau_\omega} V(s')\right]\). Esto nos permite aplicar Q-learning (o Actor-Critic) en la política de alto nivel sin tener que hacer actualizaciones combinatorias para cada acción atómica.

\begin{knowledgebox}{Desarrollo de Bellman semimarkoviano}
\(\quad V(s)=\sum_\omega \pi_{\text{high}}(\omega|s)\left[\mathbb{E}_{\tau_\omega}\left[\sum_{t=0}^{\tau_\omega}\gamma^t r_t\right]+\mathbb{E}[\gamma^{\tau_\omega}V(s')]\right]\)
\end{knowledgebox}

\begin{warningbox}{Sobreajuste de las options}
Si $\mathcal{I}_\omega$ es demasiado estrecho o la estructura de $\pi_\omega$ es demasiado rígida, la política de alto nivel cae en local minima del tipo «solo unas cuantas options son eficaces». Se recomienda aplicar regularization (dropout style) o ampliar dinámicamente las options disponibles.
\end{warningbox}

\begin{table}[H]
\centering
\begin{tabular}{p{4cm}p{5cm}}
\toprule
Componente & Función \\
\midrule
Initiation set & Controla cuándo puede el nivel alto elegir una option \\
Intra-option policy & Se encarga de la ejecución de varios pasos, junto con el skill replay \\
Termination function & Decide si la option se mantiene \\
High-level policy & Elige options con una granularidad temporal gruesa \\
\bottomrule
\end{tabular}
\caption{Componentes clave del marco de options.}
\end{table}

\subsection{Mantenimiento y ampliación del conjunto de options}

En proyectos grandes hay que mantener un conjunto de options: el skill engineer actualiza el initiation set y la función de termination según la coverage/usage matrix, y periódicamente retira las options con baja tasa de uso o con un failure rate alto. El usage log puede escribirse en un instrumentation dashboard para facilitar la governance review.

\begin{figure}[H]
\centering
\includegraphics[width=0.85\textwidth]{slides-images/slide-08.jpg}
\caption{Flujo de entrenamiento de options y usage tracking (slide-08).}
\end{figure}

\begin{importantbox}{Recomendaciones para operar el conjunto de options}
\begin{itemize}
    \item Usar una coverage matrix para supervisar cuántas veces se invoca cada option;
    \item Cada cambio en una option debe actualizar a la vez el prompt log y la instrumentation;
    \item Las options nuevas se prueban primero en un sandbox y después se pasan a production;
\end{itemize}
\end{importantbox}

\subsection{Resumen de la sección}

Options/SMDP aporta la formalización matemática del RL jerárquico: pone el acento en la extensión temporal de las action y en la reformulación de la función value, y en los proyectos reales exige registrar initiation/termination/logs en la instrumentation.
